% Smoke-test deck for the code_to_tex Beamer + overlay system: a subset of the old
% spork/docs/spork_nyc25.tex (contents need not match nyc25.py; formatting is the point).
% Build:  cd spork/docs && python3 -m code_to_tex.build code_to_tex/beamer_nyc25_test --slides beamer_nyc25_test --png
% One frame:  ... --slides beamer_nyc25_test --frame dag --skip-code --png
\documentclass[aspectratio=169]{beamer}
\newif\ifsporkSlides \sporkSlidestrue
\input{copied_styles/spork_style.tex}
\input{copied_styles/code_to_tex_support.tex}

\title{Exo-GPU}
\subtitle{Safe, Imperative, User-schedulable Programming for Tensor Cores}
\author{David Zhao Akeley \and Yuka Ikarashi \and Jonathan Ragan-Kelley}
\date{2025 MIT/Jane Street Symposium (smoke-test port)}

\begin{document}

% ---------------------------------------------------------------------------
\begin{frame}[label=title]
\titlepage
\end{frame}

% ---------------------------------------------------------------------------
% Listing with highlight steps: version v on step v+1; version 4 removes the signature lines.
% The title changes per step (\sporkSlideOnly in the title).
\begin{frame}[label=cpu]{%
    \sporkSlideOnly<1>{Starting out with Exo}%
    \sporkSlideOnly<2>{Type Annotations}%
    \sporkSlideOnly<3>{Memory Type Annotations}%
    \sporkSlideOnly<4>{M/N Tiling}%
    \sporkSlideOnly<5>{Loop Body}%
    \sporkTitleTag{\sporkInvBlue{GEMM v1}}}
\footnotesize
\sporkListing{cpu}
\normalsize
\vspace{2mm}
\sporkCodeRule
\vspace{2mm}
\sporkSlideOnly<1>{\texttt{@proc} (procedure) decorator: Python AST $\to$ C/C++ (mostly 1:1).}%
\sporkSlideOnly<2>{Sizes and tensor shapes are \sporkKeyA{type annotations}.}%
\sporkSlideOnly<3>{\texttt{@ DRAM}: every buffer has a \sporkKeyB{memory type}.}%
\sporkSlideOnly<4>{We want to \sporkKeyA{tile} the \sporkGreen{m}/\sporkViolet{n} loop nest.}%
\sporkSlideOnly<5>{\sporkKeyA{Sequential} loops here; contrast to \sporkKeyB{parallel} loops later.}
\end{frame}

% ---------------------------------------------------------------------------
% `reserve`: the assert line appears in version 2, and nothing below it moves.
\begin{frame}[label=m_divide]{%
    \sporkSlideOnly<1>{Divide Loop (M)}%
    \sporkSlideOnly<2>{M Variable Rewritten}%
    \sporkSlideOnly<3>{Assume Perfect Tiling}}
\footnotesize
\sporkListing[reserve]{m_divide_loop}
\end{frame}

% ---------------------------------------------------------------------------
% Fixed versions of two listings side by side (no overlays).
\begin{frame}[label=before_after]{Before and After (Exo checked functional equivalence)}
\begin{columns}[T]
\begin{column}{0.4\textwidth}
\textbf{BEFORE}\par
\fontsize{9.5}{11.2}\selectfont
\sporkListing[version=4]{cpu}
\end{column}
\begin{column}{0.56\textwidth}
\textbf{AFTER}\par
\fontsize{9.5}{11.2}\selectfont
\sporkListing[version=1]{m_divide_loop}
\end{column}
\end{columns}
\end{frame}

% ---------------------------------------------------------------------------
% Listing beside a figure: listing in a column (normal text), figure in the other column.
% Both animate: listing version 1 and the figure's A/B tiles appear on step 2.
\begin{frame}[label=grid_tile]{Summary: Thread Block Tiling\sporkTitleTag{\sporkInvBlue{GEMM v2}}}
\begin{columns}[T]
\begin{column}{0.48\textwidth}
\footnotesize
\sporkListing{loops}
\end{column}
\begin{column}{0.48\textwidth}
\centering
\scalebox{0.95}{% C tiled into M1 x N1 thread-block tiles (from the old nyc25_grid_tile.tex + the
% "Block Matrix Product" slide). Overlay steps (frame `grid_tile`):
%   1  C with the block grid                                     (paper: shown)
%   2+ A row tiles and B column tiles   spork reveal             (paper: shown)
%   2  block 0 emphasized               spork slide only=<2>{...} (paper: plain)
\documentclass[beamer,aspectratio=169]{standalone}
\newif\ifsporkSlides \sporkSlidestrue
\input{copied_styles/spork_style.tex}
\input{copied_styles/code_to_tex_support.tex}
\sporkFigMaxWidth{7cm}
\begin{document}
\begin{standaloneframe}
\begin{tikzpicture}[
    spork figure,
    cell/.style={minimum size=10mm, inner sep=0pt, anchor=north west},
  ]
  % C: 4 x 4 grid of M1 x N1 block tiles; (0,0) is the top left corner.
  \draw [thick, fill=white] (0, 0) rectangle (40mm, -40mm);
  \draw [gray] (0, 0) grid[step=10mm] (40mm, -40mm);
  \node at (20mm, -25mm) {\fontsize{20}{24}\selectfont C};
  \node [cell, spork slide only=<2>{spork inv green style, draw}] at (0, 0) {0};
  \node [cell] at (10mm, 0) {1};
  \node [cell] at (20mm, 0) {2};
  \node [cell] at (30mm, 0) {\dots};
  \node [cell] at (0, -10mm) {\dots};
  % Tile size labels.
  \node [anchor=west] at (40mm, -5mm) {\sporkGreen{\texttt{M1}}};
  \node [anchor=north] at (5mm, -40mm) {\sporkViolet{\texttt{N1}}};
  \node [anchor=west] at (40mm, -25mm) {\sporkGreen{\texttt{M}}};
  \node [anchor=north] at (25mm, -40mm) {\sporkViolet{\texttt{N}}};
  % Step 2+: the A and B tiles each block row/column reads.
  \foreach \r in {0,...,3} {
    \node [spork reveal=<2->, spork green style, draw, minimum width=14mm, minimum height=9mm, anchor=east]
      at (-2mm, -5mm - \r*10mm) {A[\dots]};
    \node [spork reveal=<2->, spork violet style, draw, minimum width=9mm, minimum height=14mm, anchor=south]
      at (5mm + \r*10mm, 2mm) {\rotatebox{90}{B[\dots]}};
  }
\end{tikzpicture}
\end{standaloneframe}
\end{document}
}
\end{column}
\end{columns}
\vfill
\sporkReveal<2->{Block \sporkInvGreen{0} reads a row of \sporkGreen{A} tiles and a column of \sporkViolet{B} tiles.}
\vspace{3mm}
\end{frame}

% ---------------------------------------------------------------------------
% Bullet text: keys, boxes, revealed items.
\begin{frame}[label=smem]{Shared Memory: SMEM}
\sporkKeyB{SMEM:} per-thread-block ``scratchpad''.

Save the needed block of \sporkGreen{A}, \sporkViolet{B} in SMEM.
\begin{itemize}
  \item Break blocks into tiles to fit within SMEM.
  \item \sporkReveal<2->{Tile the \sporkBlue{\texttt{k}} loop by \sporkBlue{\texttt{K0}}
        (\texttt{k $\mapsto$ k1 * K0 + k0}).}
  \item \sporkReveal<3->{Copies: \sporkKeyA{GMEM} $\to$ \sporkKeyB{SMEM}, then many reads from SMEM.}
  \begin{itemize}
    \item \sporkReveal<3->{\sporkYellow{\texttt{Fence}} needed: the thread that \emph{writes} a value is not the one that \emph{reads} it.}
  \end{itemize}
\end{itemize}
\vfill
\sporkReveal<4->{\sporkMainKey{Main key}, \sporkMinorKey{minor key}; boxes:
  \sporkRed{red} \sporkYellow{yellow} \sporkGreen{green} \sporkBlue{blue} \sporkViolet{violet}
  \sporkInvRed{inv red} \sporkInvGold{inv gold} \sporkInvGreen{inv green} \sporkInvPurple{inv purple}.}
\vspace{3mm}
\end{frame}

% ---------------------------------------------------------------------------
% Code grows: the Fence lines appear on step 2 (elided arguments), then highlights.
\begin{frame}[label=smem_sync]{%
    \sporkSlideOnly<1>{Stage in SMEM: Divide \& Parallelize SMEM-fill Loops}%
    \sporkSlideOnly<2-3>{Add Synchronization\sporkTitleTag{\sporkInvBlue{GEMM v3}}}%
    \sporkSlideOnly<4->{Synchronization: What Is Ordered?\sporkTitleTag{\sporkInvBlue{GEMM v3}}}}
\footnotesize
\sporkListing{smem_in_order}
\end{frame}

% ---------------------------------------------------------------------------
\begin{frame}[label=sync_flow]{%
    \sporkSlideOnly<1>{Loop Fission Restructuring}%
    \sporkSlideOnly<2>{``Classic'' GPU GEMM Summary}%
    \sporkSlideOnly<3>{Modernize GPU GEMM}}
% Fixed-height text area: alternatives of different heights would make the figure jump.
\begin{overlayarea}{\linewidth}{2\baselineskip}
\sporkSlideOnly<1>{All steps are thread-block-cooperative, but nothing orders them yet.}%
\sporkSlideOnly<2>{First sync: wait for the SMEM tile to fill (RAW).\\
  Second sync: don't overwrite until all SMEM reads are done (WAR).}%
\sporkSlideOnly<3>{Next step: use accelerator instructions for the \sporkKeyA{GMEM}$\to$\sporkKeyB{SMEM} copies.}
\end{overlayarea}
\begin{center}
\scalebox{1.2}{% "Classic" GPU GEMM main loop as a flow (from the old nyc25_sync_flow.tex + nyc25_sync_caption.tex
% + the "Loop Fission Restructuring"/"Modernize GPU GEMM" annotations). Overlay steps:
%   1  flow without synchronization (direct B -> accumulate edge)           (slide only)
%   2+ RAW/WAR sync bars + captions        spork reveal                     (paper: shown)
%   3  "Accelerate me!" callouts on the two loads                           (slide only)
\documentclass[beamer,aspectratio=169]{standalone}
\newif\ifsporkSlides \sporkSlidestrue
\input{copied_styles/spork_style.tex}
\input{copied_styles/code_to_tex_support.tex}
\sporkFigMaxWidth{11cm}
\begin{document}
\begin{standaloneframe}
\begin{tikzpicture}[
    spork figure,
    node distance=5mm,
    flow box/.style={minimum width=22mm, minimum height=17mm, align=center},
    flow sync/.style={spork yellow style, minimum width=4mm, minimum height=17mm, spork reveal=<2->},
  ]
  \node (A) [flow box, spork green style] {Load A tile\\to SMEM\\(\texttt{M1 $\times$ K0})};
  \node (B) [flow box, spork violet style, right=of A] {Load B tile\\to SMEM\\(\texttt{K0 $\times$ N1})};
  \node (raw) [flow sync, right=of B] {};
  \node (accum) [flow box, draw=black, right=of raw] {Accumulate\\tile\\(\sporkBlue{\texttt{k0}} loop)};
  \node (war) [flow sync, right=of accum] {};

  \node (zero) [above=8mm of A] {Zero accumulators};
  \node (C) [above=8mm of war] {Write out C};

  \draw [spork arrow] (zero) -- (A);
  \draw [spork arrow] (war) -- (C);
  \draw [spork arrow] (A) -- (B);
  \draw [spork arrow, spork reveal=<2->] (B) -- (raw);
  \draw [spork arrow, spork reveal=<2->] (raw) -- (accum);
  \draw [spork arrow, spork slide only=<1>] (B) -- (accum);
  \draw [spork arrow, spork reveal=<2->] (accum) -- (war);
  \draw [spork arrow] (war.north) to[out=150, in=30]
    node[midway, above] {\sporkBlue{\texttt{k1}} loop} (A.north);

  \node (rawCaption) [below=1mm of raw, align=center, spork reveal=<2->, font=\bfseries]
    {Sync:\\RAW};
  \node (warCaption) [below=1mm of war, align=center, spork reveal=<2->, font=\bfseries]
    {Sync:\\WAR};
  \node [below=1mm of $(A.south)!0.5!(B.south)$, align=center, font=\bfseries]
    {All steps are thread-\\block-cooperative.};

  % Step 3 (slide only): what the next part of the talk accelerates.
  \foreach \n in {A, B} {
    \node (\n accel) [spork slide only=<3>, spork red style, above=3mm of \n, xshift=4mm,
      font=\bfseries] {Accelerate me!};
  }
\end{tikzpicture}
\end{standaloneframe}
\end{document}
}
\end{center}
\end{frame}

% ---------------------------------------------------------------------------
% The merged DAG figure: steps 1-3 build the paper figure, 4-5 are slide-only emphasis.
\begin{frame}[label=dag]{%
    \sporkSlideOnly<1>{Loop Skew: Delay SMEM Tile Consumption}%
    \sporkSlideOnly<2>{Ring Buffer Required}%
    \sporkSlideOnly<3>{WAR Hazard: Ring Buffer Slot Re-Use}%
    \sporkSlideOnly<4>{Split Barriers}%
    \sporkSlideOnly<5>{Tensor Cores}}
\begin{center}
\scalebox{1.2}{% Loop-skewed producer/consumer DAG (merged from the old nyc25_dag.tex, nyc25_dag_k.tex,
% nyc25_dag_ring.tex, nyc25_dag_war.tex). Overlay steps (frame `dag` of the deck):
%   1  loads, computes, RAW edges, k1 labels          (paper: shown)
%   2+ ring buffer slot row        spork reveal       (paper: shown)
%   3+ WAR edges                   spork reveal       (paper: shown, dashed via spork paper style)
%   4  one RAW + one WAR edge and their nodes emphasized, "split barrier" callout   (slide only)
%   5  compute nodes red, "Accelerate me!" callout                                   (slide only)
% Paper only: a legend explaining solid/dashed edges.
\documentclass[beamer,aspectratio=169]{standalone}
\newif\ifsporkSlides \sporkSlidestrue
\input{copied_styles/spork_style.tex}
\input{copied_styles/code_to_tex_support.tex}
\sporkFigMaxWidth{11cm}
\begin{document}
\begin{standaloneframe}
% `spork figure`: absolute font size (\sporkFigFont), the same in the paper (10pt body) and on slides (12pt body).
\begin{tikzpicture}[
    spork figure,
    dag load/.style={minimum width=15mm, minimum height=5.5mm, inner sep=1pt},
    dag compute/.style={draw=black, fill=white, minimum width=15mm, minimum height=10mm},
    % Per-node one-off emphasis (slide only); referenced as `dag X/.try` below.
    dag C2/.style={spork slide only=<4>{spork red style, very thick}},
    dag B4/.style={spork slide only=<4>{spork inv purple style}},
    dag A1/.style={spork slide only=<4>{spork inv green style}},
  ]
  \def\pitch{17mm}
  % Loads (iteration i) and computes (iteration i, consuming the tile loaded in i-1).
  \foreach \i in {0,...,5} {
    \node (A\i) [dag load, spork green style, dag A\i/.try] at (\i*\pitch, 0) {Load A};
    \node (B\i) [dag load, spork violet style, dag B\i/.try, anchor=north] at (A\i.south) {Load B};
    \node (k\i) at (\i*\pitch, -27mm) {\sporkBlue{\texttt{k1 = \i}}};
  }
  \foreach \i in {1,...,5} {
    \node (C\i) [dag compute, spork slide only=<5>{spork red style}, dag C\i/.try] at (\i*\pitch, -19mm) {compute};
  }
  % RAW edges: tile loaded in iteration i is consumed in iteration i+1.
  \foreach \i [evaluate=\i as \j using int(\i+1)] in {0,...,4} {
    \draw [spork arrow] (A\i.south east) to[out=0, in=180] (C\j.west);
  }
  % Ring buffer slots (3 slots; slot reuse boxed).
  \node (buf) [spork reveal=<2->, align=right, font=\bfseries, anchor=east]
    at ($(k0.east)+(0,-6mm)$) {Ring\\buffer:};
  \foreach \i/\s/\boxed in {1/0/1, 2/1/0, 3/2/0, 4/0/1, 5/1/0} {
    \ifnum\boxed=1
      \node (slot\i) [spork reveal=<2->, draw=black, inner sep=2pt] at ($(k\i)+(0,-6mm)$) {Slot \s};
    \else
      \node (slot\i) [spork reveal=<2->, inner sep=2pt] at ($(k\i)+(0,-6mm)$) {Slot \s};
    \fi
  }
  % WAR edges: compute i must finish before slot (i mod 3) is refilled in iteration i+2.
  \foreach \i [evaluate=\i as \j using int(\i+2)] in {1,...,3} {
    \draw [spork arrow, spork reveal=<3->, spork paper style={dashed}]
      ($(C\i.north east)+(0,-2mm)$) to[out=40, in=200] (B\j.west);
  }
  % Step 4 (slide only): one RAW and one WAR edge redrawn thick red, plus a callout.
  \draw [spork arrow, very thick, draw=sporkRedFg, spork slide only=<4>]
    (A1.south east) to[out=0, in=180] (C2.west);
  \draw [spork arrow, very thick, draw=sporkRedFg, spork slide only=<4>]
    ($(C2.north east)+(0,-2mm)$) to[out=40, in=200] (B4.west);
  \node (splitNote) [spork slide only=<4>, spork red style, draw, align=center, anchor=north east]
    at ($(slot5.south east)+(0,-1mm)$) {each edge = one\\\texttt{Arrive}/\texttt{Await} pair};
  % Step 5 (slide only): tensor core callout.
  \node [spork slide only=<5>, spork inv red style, font=\bfseries]
    at ($(C3.north)!0.5!(C4.north)+(0,3mm)$) {Accelerate me!};
  % Paper only: legend (on slides the speaker explains it).
  \node [spork paper only, anchor=north west, align=left]
    at ($(slot1.south west)+(0,-2mm)$)
    {Solid edges: RAW (tile ready); dashed edges: WAR (ring buffer slot free).};
\end{tikzpicture}
\end{standaloneframe}
\end{document}
}
\end{center}
\vfill
\sporkCodeRule
\vspace{2mm}
\sporkSlideOnly<1>{Consume each SMEM tile one \sporkBlue{\texttt{k1}} iteration \sporkKeyA{after} it is loaded.}%
\sporkSlideOnly<2>{One slot may be being filled while another is consumed.}%
\sporkSlideOnly<3>{A slot can only be refilled once its previous readers are done.}%
\sporkSlideOnly<4>{Implement each edge with a \sporkKeyA{split barrier}: \texttt{Arrive} \dots{} \texttt{Await}.}%
\sporkSlideOnly<5>{Use tensor cores to accelerate \sporkKeyA{matrix multiply-accumulate}.}
\vspace{3mm}
\end{frame}

% ---------------------------------------------------------------------------
% Listing marks + arrows between code lines; remark lines (plain and bold); placeholders.
\begin{frame}[label=split_marks]{Split Barriers in Code}
% The `split` listing with arrows drawn between its Arrive and Await lines (listing marks;
% same technique as the arXiv paper's async-bug example, without the \bugYFix hack).
% A "figure" that holds a listing: \input by the deck (frame `split_marks`) and by main.tex.
%   slides: listing version v on step v+1; RAW arrow from step 2, WAR arrow from step 3
%   paper:  last listing version, both arrows
% Marks must be in normal text (not inside a TikZ node), so the listing is plain text and
% the arrows live in a separate `remember picture, overlay` picture (code_to_tex/README.md, Marks).
% (Like every figure, the standalone preview uses the deck's 16:9 page; standalone passes the
% option on to Beamer. Nothing may precede \documentclass, since that would run when the figure is \input.)
\documentclass[beamer,aspectratio=169]{standalone}
\newif\ifsporkSlides \sporkSlidestrue
\input{copied_styles/spork_style.tex}
\input{copied_styles/code_to_tex_support.tex}
\begin{document}
\begin{standaloneframe}
\begin{minipage}{116mm}
% Absolute size: \footnotesize differs between the paper (8pt) and slides (10pt).
\fontsize{9}{10.5}\selectfont
\sporkListing{split}
\end{minipage}%
\begin{tikzpicture}[remember picture, overlay, rounded corners=2mm,
    spork figure, font=\bfseries]
  % Vertical arrow segments run at fixed distances from the listing's left edge (mark `left`),
  % right of the longest code line they pass.
  \coordinate (rawX) at ([xshift=88mm]pic cs:split-left);
  \coordinate (warX) at ([xshift=116mm]pic cs:split-left);
  \foreach \m in {raw_arrive, raw_await, war_arrive, war_await} {
    \coordinate (\m) at ([xshift=1mm]pic cs:split-\m);
  }
  \sporkReveal<2->{
    \draw [spork arrow, draw=sporkRedFg] (raw_arrive) -- (raw_arrive -| rawX)
      -- node[right, text=sporkRedFg] {RAW} (raw_await -| rawX) -- (raw_await);
  }
  \sporkReveal<3->{
    \draw [spork arrow, draw=sporkBlueFg] (war_arrive) -- (war_arrive -| warX)
      -- node[right, text=sporkBlueFg, align=left] {WAR\\(slot\\reuse)} (war_await -| warX) -- (war_await);
  }
\end{tikzpicture}
\end{standaloneframe}
\end{document}

\end{frame}

% ---------------------------------------------------------------------------
\sporkExtraSlides

\begin{frame}[label=extra_ring]{Extra Slides: %
    \sporkSlideOnly<1>{Ring Buffer}\sporkSlideOnly<2>{Loop Skew}}
\scriptsize
\sporkListing{ring}
\normalsize
\vfill
\sporkSlideOnly<1>{Lift \texttt{A\_smem}, \texttt{B\_smem} out of the \sporkBlue{\texttt{k1}} loop and ring buffer by 3
  (long index expressions \sporkComment{elided}).}%
\sporkSlideOnly<2>{Add an extra \sporkBlue{\texttt{k1}} iteration; delay the accumulate code by one iteration.}
\vspace{3mm}
\end{frame}

\begin{frame}[label=extra_machine]{Extra Slides: Abstract Machine}
Exo-GPU programs target this abstract machine.
\medskip

\textbf{Correctness view:} check predicates on \sporkKeyA{visibility sets}.
\begin{itemize}
  \item Each read is \sporkKeyA{guaranteed} to get the value of the sequentially-prior write.
  \item Sequential-parallel equivalence.
\end{itemize}
\medskip

\textbf{Optimization view:} the abstract machine trace \sporkKeyB{dictates} which instructions and threads perform each operation.
\begin{itemize}
  \item Constrains the compiler.
  \item Generated CUDA internals aren't implementation details.
\end{itemize}
\end{frame}

\end{document}
